\documentclass[11pt]{article}
\usepackage[margin=1in]{geometry}
\usepackage{amsmath,amssymb,amsthm}
\usepackage{hyperref}
\usepackage{listings}
\usepackage{xcolor}
\usepackage{booktabs}
\usepackage{array}
\usepackage{microtype}
\usepackage{cite}
\usepackage{url}
\usepackage{fancyhdr}
\usepackage{skak}        % chess symbols
\usepackage{chessboard}  % board diagrams from FEN

% Allow line breaks at underscores and dots in \texttt
\renewcommand{\texttt}[1]{%
  \begingroup\ttfamily\hyphenchar\font=`\-\relax
  \def\_{\textunderscore\discretionary{}{}{}}%
  #1\endgroup}

% Global stretch to prevent overfull boxes from long identifiers
\emergencystretch=3em

\newtheorem{theorem}{Theorem}
\newtheorem{lemma}[theorem]{Lemma}
\newtheorem{proposition}[theorem]{Proposition}
\newtheorem{corollary}[theorem]{Corollary}
\theoremstyle{definition}
\newtheorem{definition}[theorem]{Definition}
\newtheorem{finding}[theorem]{Finding}
\theoremstyle{remark}
\newtheorem{remark}[theorem]{Remark}

\lstset{
  basicstyle=\ttfamily\small,
  keywordstyle=\color{blue},
  commentstyle=\color{gray},
  breaklines=true,
  frame=single,
  xleftmargin=1em,
  escapeinside={(*}{*)},
}

\title{Yanasse: Finding New Proofs from Deep Vision's Analogies\\[0.5em]
\large Part~3: Probability Theory $\to$ Condensed Mathematics}
\author{Alexandre Linhares\thanks{alexandre@linhares.ltd. Affiliations: Getulio Vargas Foundation and ARGO LABS}}
\date{April 2026}

\pagestyle{fancy}
\fancyhf{}
\renewcommand{\headrulewidth}{0pt}
\fancyfoot[C]{\thepage}
\fancyfoot[L]{\footnotesize\textit{Citation: Linhares, A. ``Yanasse: Finding New Proofs from Deep Vision's Analogies --- Part~3,'' Argolab.org Report for Dissemination, 2026.}}

\begin{document}
\maketitle

\begin{abstract}
This is Part~3 of the Yanasse series, applying cross-area tactic
transfer to the pair \textsc{Probability}~$\to$~\textsc{Condensed}.  Condensed mathematics---the sheaf-theoretic
framework of Clausen and Scholze that replaces point-set
topology---is structurally the most distant target tested so far:
it was designed to \emph{eliminate} the filter-based and
$\sigma$-algebra-based vocabulary on which probability-theory
tactics depend.

The method produces 2~Lean-verified new proofs out of 10 transfer
attempts (20\%), the lowest success rate in the series.  Both
proofs target the same theorem---the triangle identity of the
adjunction between condensed sets and topological spaces
(\texttt{topCatAdjunction.left\_triangle\_components})---and both
transfer the ``go pointwise'' pattern that probability theory uses
to reduce equalities of structured objects (kernels, random
variables) to per-element obligations.

The central finding is a \emph{strictly better proof}: the
transferred tactic \texttt{congr with x} replaces \texttt{ext}
\emph{without} requiring the
\texttt{set\_option backward.isDefEq.respectTransparency false}
workaround that the existing Mathlib proof needs.  This is not a
paraphrase of the original proof but a genuine improvement enabled
by cross-area transfer.  The remaining 8~schemas all fail for a
common root cause: condensed mathematics systematically avoids the
type-theoretic structures (filters, $\sigma$-algebras, conditional
definitions, \texttt{CanLift} instances) that probability-theory
tactics operate on.

Cross-pair comparison across all three target areas confirms a
sharp dichotomy: schemas that interact with measure-theoretic
\emph{types} never transfer; schemas that operate on
\emph{proof architecture} (decomposition, case-split, congruence)
transfer when the target has compatible goal shapes.
\end{abstract}

\newpage
\tableofcontents
\newpage

%======================================================================
\section{Introduction}\label{sec:intro}
%======================================================================

Parts~1 and~2 of this series~\cite{yanasse_part1, yanasse_part2}
demonstrated cross-area tactic transfer at a 40\% schema success
rate on Probability $\to$ Representation Theory and Probability
$\to$ Set Theory alike.  Both target areas, while
mathematically distant from probability, share a common property:
they use the standard vocabulary of Lean's type theory---equalities,
inequalities, type-class instances---in ways that probability-theory
proof patterns can engage with.

This paper tests the method against a fundamentally different
target.  Condensed mathematics, developed by Clausen and
Scholze~\cite{scholze_condensed_2019, clausen_scholze_complex_2022},
replaces point-set topology with a sheaf-theoretic framework: a
condensed set is a sheaf on the category of compact Hausdorff
spaces.  The theory was designed precisely to avoid the categorical
deficiencies of topological spaces (the category $\mathbf{Top}$
lacks exactness properties), and in doing so it eliminates the
filter-based and measure-theoretic vocabulary that probability-theory
tactics depend on.  This makes \textsc{Condensed} the hardest
target in the Yanasse series: every domain-gated schema faces not
just an unfamiliar target but an \emph{actively hostile} one.

The Lean formalization of condensed mathematics in
Mathlib~\cite{asgeirsson_condensed_2024, asgeirsson_discrete_2024}
is comparatively young: 189~extracted proof states across 33~files,
versus 672 for Representation Theory and 2{,}076 for Set Theory.
This small pool increases the risk of QAP score saturation (multiple
targets tied at the same score) and limits the diversity of available
target theorems.

\subsection{Overview of the Deep Vision framework}

The Deep Vision framework~\cite{linhares_deep_vision,
linhares2024deepvision}, growing out of the Fluid Analogies Research
Group tradition~\cite{hofstadter_fluid_1995, hofstadter_go_1999,
hofstadter_surfaces_2013}, represents any structured domain as a
\emph{relational network}: entities connected by typed, directed
relations.  An analogy optimization finds the entity permutation that
maximizes relational consistency between two networks.  The matching
code (\texttt{deep\_vision\_lib.py}) is entirely domain-independent:
it receives two relational structures and returns a similarity score
and entity correspondence, without knowing which domain it is
processing.

\paragraph{Chess analogies.}  In chess~\cite{linhares_active_2005,
linhares_understanding_2007, linhares_questioning_2010,
linhares_entanglement_2012, linhares_what_2013,
linhares_emergence_2014}, entities are pieces and relations are attack,
defense, blocking, confinement, and pinning.  Two pieces in entirely
different positions can play ``the same role'' if their relational
profiles match.

\paragraph{From chess to proofs.}  The same matching engine, with only
the relation extractor swapped, handles Lean~4 proof states: entities
become hypotheses, goals, and types; relations become rewrite,
fit/apply, head-match, structure, equality, reflexive, witness,
bidirectional, kernel/simp, decidable, and lemma-needed (14 relation
types).  The proof-state matcher has previously been
used to discover a formal proof of Wolstenholme's
theorem~\cite{linhares_wolstenholme_2026} in Lean~4
(see also~\cite{linhares_sketch_2026}).

\subsection{Results preview}

Applying the method to Probability $\to$ Condensed produces:

\begin{itemize}
\item 2 Lean-verified new proofs from 2 successful schema transfers
  out of 10 attempts (20\% transfer success rate).
\item A \emph{strictly better proof}: the transferred
  \texttt{congr with x} avoids a transparency workaround
  (\texttt{set\_option}) that the existing Mathlib proof requires
  (Section~\ref{sec:proof1}).
\item A paradigm-incompatibility diagnosis: 8 of 10 schemas fail
  because condensed mathematics systematically eliminates the
  type-theoretic structures they require
  (Section~\ref{sec:failures}).
\item Three-way cross-pair comparison across all target areas tested
  so far (Section~\ref{sec:cross_pair}).
\end{itemize}

%======================================================================
\section{Methodology}\label{sec:method}
%======================================================================

The methodology is identical to Parts~1 and~2~\cite{yanasse_part1,
yanasse_part2}.  We summarize the pipeline and note the parameters
specific to this pair.

\subsection{Pipeline summary}

\begin{enumerate}
\item \textbf{Schema extraction.}  Parse the Mathlib proof-state
  corpus (217{,}133 entries).  Each tactic invocation becomes a schema
  tuple $(\textit{head}, \textit{arity}, \textit{has\_with},
  \textit{uses\_lemma})$.

\item \textbf{$z$-score ranking.}  For each (area, schema) pair,
  compute the $z$-score of the schema's frequency across 27 areas.
  A schema with $z \geq +2$ in the source and $z \leq -1$ (or absent)
  in the target is a transfer candidate.

\item \textbf{GPU-accelerated analogy matching.}  Sample 8 source
  proofs per schema, match each against all 189 Condensed proof
  states using the \texttt{BatchMatcher} on MPS.  Retain the top-5
  target analogues by normalized score.

\item \textbf{Semantic adaptation.}  A fresh AI reasoning session
  (Claude Opus, 50-minute budget) adapts the source tactic's
  \emph{mathematical strategy}---not its syntax---to the target
  theorem.

\item \textbf{Lean verification.}  The adapted proof is compiled
  with \texttt{lake env lean}.  Exit code~0 with no \texttt{sorry}
  constitutes verification.
\end{enumerate}

\subsection{Pair-specific parameters}

\begin{itemize}
\item \textbf{Source:} \texttt{Mathlib.Probability} (30{,}384 tactic
  invocations, the largest Mathlib area).
\item \textbf{Target:} \texttt{Mathlib.Condensed} (189 proof
  states---the smallest target pool in the series, ${\sim}9\times$
  smaller than Set Theory and ${\sim}3.5\times$ smaller than
  Representation Theory).
\item \textbf{Matching cost:} $8 \times 189 = 1{,}512$
  individual analogy computations per schema, completing in
  approximately 2--10~minutes on MPS---significantly faster
  than the previous pairs due to the smaller pool.
\item \textbf{Top-10 transfer candidates} (by gap score):
  \texttt{filter\_upwards}, \texttt{congr\,with}, \texttt{fun\_prop}
  (as \texttt{all\_goals fun\_prop}), \texttt{ext1},
  \texttt{lift}, \texttt{any\_goals},
  \texttt{measurability}, \texttt{by\_cases}, \texttt{congrm},
  \texttt{split\_ifs}.
\end{itemize}

Note the substitutions relative to previous pairs:
\texttt{all\_goals fun\_prop} replaces bare \texttt{fun\_prop}
(the full combinator, not just the inner tactic, is the
actual source pattern), and \texttt{split\_ifs} replaces
\texttt{push} (the next-highest-gap schema for this target).

%======================================================================
\section{The Target: Condensed Mathematics}\label{sec:target}
%======================================================================

A \emph{condensed set} is a sheaf of sets on the category
$\mathbf{CompHaus}$ of compact Hausdorff spaces equipped with the
coherent topology~\cite{scholze_condensed_2019}.  The functor
\[
  T \colon \mathbf{Top} \longrightarrow \mathbf{CondSet},
  \qquad
  X \longmapsto \bigl(S \mapsto \mathrm{Cont}(S, X)\bigr)
\]
embeds topological spaces into condensed sets, and admits a left
adjoint
\[
  U \colon \mathbf{CondSet} \longrightarrow \mathbf{Top}
\]
that recovers the ``underlying topological space'' of a condensed
set.  The adjunction $U \dashv T$ is the categorical backbone of the
Mathlib formalization~\cite{asgeirsson_condensed_2024}: its unit
$\eta \colon \mathrm{Id} \Rightarrow T \circ U$ and counit
$\varepsilon \colon U \circ T \Rightarrow \mathrm{Id}$ satisfy the
triangle identities
\begin{equation}\label{eq:triangle}
  \varepsilon_{T(X)} \circ T(\eta_X) = \mathrm{id}_{T(X)},
  \qquad
  U(\varepsilon_Y) \circ \eta_{U(Y)} = \mathrm{id}_{U(Y)}.
\end{equation}

Both new proofs in this paper target the second identity
(Equation~\ref{eq:triangle}, right), formalized in Mathlib as
\texttt{topCatAdjunction.left\_triangle\_components}.

\paragraph{Why condensed math is hard to reach from probability.}
The Clausen--Scholze framework was designed to replace filter-based
topology---the very vocabulary that probability theory's most
characteristic tactics (\texttt{filter\_upwards},
\texttt{measurability}, \texttt{fun\_prop}) operate on---with
sheaf-theoretic conditions.  Condensed proofs reason about
\emph{functors}, \emph{natural transformations}, and
\emph{adjunctions}, not about $\sigma$-algebras, filters, or
conditional definitions.  This makes the pair a stress test for
cross-area transfer: any success must come from proof patterns
that transcend the measure-theoretic vocabulary.

The Mathlib formalization draws on the work of
Asgeirsson~\cite{asgeirsson_condensed_2024, asgeirsson_discrete_2024},
building on the foundations laid by the Liquid Tensor
Experiment~\cite{scholze_liquid_2020, commelin_lte_2022}, which
formalized Scholze's theorem on liquid vector spaces in Lean and
demonstrated that condensed mathematics could be made fully
rigorous in a proof assistant.

%======================================================================
\section{Results: Probability $\to$ Condensed}\label{sec:results}
%======================================================================

\begin{table}[h]
\centering\small
\caption{Transfer results for Probability $\to$ Condensed.}
\label{tab:results}
\begin{tabular}{@{}clcccp{4.5cm}@{}}
\toprule
Item & Schema head & \texttt{with} & Ar.\ & Result & Adaptation \\
\midrule
20 & \texttt{filter\_upwards} & \checkmark & 1 & Failed & Paradigm gap (no Filter API) \\
21 & \texttt{congr} & \checkmark & 0 & \textbf{New proof} & \texttt{congr with x} (no \texttt{set\_option}) \\
22 & \texttt{ext1} & & 1 & \textbf{New proof} & \texttt{ext1 x} (named variable) \\
23 & \texttt{all\_goals} & & 1 & Failed & No multi-goal sites \\
24 & \texttt{lift} & \checkmark & 5 & Failed & No \texttt{CanLift} instances \\
25 & \texttt{any\_goals} & & 1 & Failed & Single sequential goal \\
26 & \texttt{measurability} & & 0 & Failed & No $\sigma$-algebra predicates \\
27 & \texttt{by\_cases} & & 4 & Failed & No conditional definitions \\
28 & \texttt{congrm} & & 5 & Failed & No shared outer operators \\
29 & \texttt{split\_ifs} & & 0 & Failed & Zero \texttt{ite}/\texttt{dite} \\
\bottomrule
\end{tabular}
\end{table}

\noindent Two transfer attempts succeeded, both targeting the same
theorem: the left triangle identity of the condensed-set/topological-space
adjunction (\texttt{topCatAdjunction.left\_triangle\_components}).

\subsection{New proof P1: \texttt{congr with x} replaces \texttt{ext}
  without transparency workaround}\label{sec:proof1}

\paragraph{Theorem.}
For a condensed set $Y$, the triangle identity
\[
  U(\varepsilon_Y) \circ \eta_{U(Y)} = \mathrm{id}_{U(Y)}
\]
holds in $\mathbf{Top}$, where $U$ = \texttt{condensedSetToTopCat},
$\eta$ = \texttt{topCatAdjunctionUnit}, and
$\varepsilon$ = \texttt{topCatAdjunctionCounit}.

\paragraph{Context.}
The triangle identities are the defining equations of an
adjunction~\cite{mac_lane_categories_1998}.  In the condensed
setting, they encode the fact that the ``condensify then
topologize'' round-trip is coherent: applying the unit and then
the counit (or vice versa) recovers the identity.  The Mathlib
formalization follows
Asgeirsson~\cite{asgeirsson_condensed_2024, asgeirsson_discrete_2024},
who developed the condensed-set infrastructure building on the
Liquid Tensor Experiment~\cite{scholze_liquid_2020, commelin_lte_2022}.
The specific theorem \texttt{left\_triangle\_components} appears in
\texttt{Mathlib.Condensed.TopCatAdjunction}, one of the core files
establishing the $U \dashv T$ adjunction.

\paragraph{The existing Mathlib proof.}
The current proof requires a global transparency override:
\begin{lstlisting}
set_option backward.isDefEq.respectTransparency false in
theorem left_triangle_components :
    condensedSetToTopCat.map (topCatAdjunctionUnit Y) (*$\gg$*)
      topCatAdjunctionCounit Y.toTopCat = (*$\mathbb{1}$*) _ := by
  ext
  change Y.obj.map ((*$\mathbb{1}$*) _) x = x
  simp
\end{lstlisting}
Without \texttt{set\_option backward.isDefEq.respectTransparency
false}, the \texttt{ext}-based proof fails: \texttt{ext} applies
the registered \texttt{@[ext]} lemma for \texttt{TopCat.Hom},
which leaves the goal in a form that interacts poorly with Lean's
backward definitional-equality transparency settings.

\paragraph{New proof (congr with x).}
The transferred tactic replaces \texttt{ext} with
\texttt{congr with x}, which operates via congruence closure
rather than extensionality lemma application:
\begin{lstlisting}
-- No set_option needed
example (Y : CondensedSet.{u}) :
    condensedSetToTopCat.map (topCatAdjunctionUnit Y) (*$\gg$*)
      topCatAdjunctionCounit Y.toTopCat = (*$\mathbb{1}$*) _ := by
  congr with x
  change Y.obj.map ((*$\mathbb{1}$*) _) x = x
  simp
\end{lstlisting}

\paragraph{Why the transfer works.}
\texttt{congr with x} performs two operations: \texttt{congr}
recognizes that both sides of the morphism equality share
the same outermost constructor (\texttt{ContinuousMap.mk}---a
function paired with a continuity proof) and decomposes accordingly;
\texttt{with x} introduces a point $x$ in the underlying topological
space.  The residual goal $Y.\mathtt{obj}.\mathtt{map}\,(\mathbb{1}\
\_)\,x = x$ follows from the sheaf functoriality axiom via
\texttt{simp}.

The key difference from \texttt{ext}: congruence closure
pattern-matches on term structure directly, bypassing the
extensionality-lemma mechanism that triggers the transparency issue.
The two tactics achieve the same mathematical decomposition
(morphism equality $\to$ pointwise equality) via different
mechanisms, and the congruence route is immune to the
backward-\texttt{isDefEq} interaction.

\paragraph{Source pattern.}  In
\texttt{Mathlib.Probability.Gaussian.HasGaussianLaw.Independence},
the tactic \texttt{congr with i} decomposes an equality of
characteristic-function products
$\prod_i \varphi(X_i)(\cdot) = \prod_i \exp(\cdots)$ to
per-factor equality on the index $i$.  The pattern---``both sides
share an outermost indexed operator; \texttt{congr} decomposes,
\texttt{with} names the index''---is the probability-theory
workhorse for equalities of kernels, expectations, and products.

\paragraph{Cross-pair comparison.}  This is the third target area
for schema~138 (\texttt{congr with}):

\medskip
\begin{center}\small
\begin{tabular}{@{}llll@{}}
\toprule
Target & Head transfers? & Modifier transfers? & Net result \\
\midrule
RepTheory (Part~1) & No ($\leq$ goal) & Yes (\texttt{rintro})
  & Adapted proof \\
SetTheory (Part~2) & No (Cardinal) & Yes (positive controls)
  & Failed (wrong rank-1) \\
\textbf{Condensed (Part~3)} & \textbf{Yes} ($=$ on ContinuousMap)
  & \textbf{Yes} (\texttt{with x})
  & \textbf{Strictly better proof} \\
\bottomrule
\end{tabular}
\end{center}
\medskip

\noindent Condensed is the only target where both the head
(\texttt{congr}) and the modifier (\texttt{with}) transfer directly.
The reason: the target goal is an equality of morphisms in a
concrete category (continuous maps between topological spaces),
which has exactly the constructor-shaped structure that
\texttt{congr} expects.

Copy-and-pasting the Lean code of the proofs above to Google's
Gemini LLM with no additional context has yielded:

\begin{quote}
  That is an excellent refactor!\ Switching from \texttt{ext} to
  \texttt{congr} is a savvy move.  While \texttt{ext} tries to apply
  every possible extensionality lemma in the database,
  \texttt{congr} (congruence) is more targeted---it basically says
  ``these two sides are the same if their components are the same.''
  The fact that you were able to remove the \texttt{set\_option}
  suggests that \texttt{congr} navigated the ``transparency'' of the
  definitions more effectively than \texttt{ext} did.  In Lean,
  \texttt{congr} is often better at ``drilling through'' categorical
  structures because it focuses on the function application rather
  than just the object type.

  \paragraph{Why this version is ``Clean.''}
  \begin{description}
    \item[No Options:] Removing \texttt{set\_option} makes the proof
      more portable and less likely to cause weird unification behavior
      in other parts of your file.
    \item[Targeted:] \texttt{congr with x} is slightly more explicit
      about what you are identifying than a generic \texttt{ext}.
    \item[Frugal:] This is a very ``short'' proof for a concept that
      is mathematically deep (the first triangle identity of an
      adjunction).
  \end{description}
\end{quote}




\subsection{New proof P2: \texttt{ext1 x} with named variable}%
\label{sec:proof2}

\paragraph{Theorem.}
Same as P1: \texttt{topCatAdjunction.left\_triangle\_components}.

\paragraph{New proof (ext1 x, three variants).}
\begin{enumerate}
\item \textbf{ext1 x + change + simp.}  Identical to the original
  proof but with \texttt{ext1 x} (named variable) instead of bare
  \texttt{ext}.  Still requires \texttt{set\_option}.
\item \textbf{ext1 x + change + simp only.}  Fully explicit lemma
  list: \texttt{FunctorToTypes.map\_id\_apply},
  \texttt{condensedSetToTopCat\_obj\_carrier}, etc.
\item \textbf{ext1 x + change + rw + rfl.}  A rewrite chain
  mimicking the probability-theory pattern: \texttt{ext1 x} $\to$
  \texttt{rw [FunctorToTypes.map\_id\_apply]} $\to$ \texttt{rfl}.
\end{enumerate}

\paragraph{Source pattern.}  In probability, \texttt{ext1\,$\omega$}
reduces set equality to pointwise membership
$\forall \omega,\; \omega \in S \leftrightarrow \omega \in T$;
\texttt{ext1\,a} reduces kernel equality to pointwise measure
equality $\forall a,\; \kappa_1(a) = \kappa_2(a)$.  The
arity-1 form (naming the variable) is standard in probability
because subsequent rewrite chains reference it.

\paragraph{Adaptation.}  \texttt{ext1 x} compiles directly because
\texttt{TopCat} morphisms have a registered \texttt{@[ext]} lemma.
The transfer is syntactic (named variable convention), not semantic.
Unlike P1, this proof still requires the \texttt{set\_option}
workaround---\texttt{ext1} applies the same extensionality lemma
as bare \texttt{ext} and triggers the same transparency interaction.

\paragraph{Domain asymmetry.}  In probability, after \texttt{ext1 x},
the goal is in \emph{rewrite-normal form}: \texttt{simp} lemmas
(\texttt{map\_apply}, \texttt{prod\_apply}) directly simplify both
sides.  In condensed math, after \texttt{ext1 x}, the goal passes
through layered categorical abstractions that are
\emph{definitionally} but not \emph{rewritably} equal to the
desired form.  The \texttt{change} tactic bridges this gap by
asking Lean's kernel to verify the definitional equality.

\paragraph{Relationship to P1.}  P1 (\texttt{congr with x}) is
strictly stronger: it achieves the same mathematical decomposition
without the \texttt{set\_option} workaround.  P2 (\texttt{ext1 x})
demonstrates that the probability naming convention transfers, but
the proof is not improved.

%======================================================================
\section{Cross-Pair Comparison: Three Target Areas}%
\label{sec:cross_pair}
%======================================================================

With three pairs of the Yanasse Project completed, we can compare schema behavior across
three structurally different target areas.  Seven schemas were tested
in all three areas (Table~\ref{tab:cross_three}).

\begin{table}[h]
\centering\small
\caption{Schema transfer results across three target areas.
  \checkmark = new proof; $\times$ = failed.}
\label{tab:cross_three}
\begin{tabular}{@{}lcccc@{}}
\toprule
Schema head & RepTheory & SetTheory & Condensed & Pattern \\
\midrule
\texttt{filter\_upwards}
  & \checkmark & $\times$ & $\times$ & Mixed \\
\texttt{congr\,with}
  & \checkmark & $\times$ & \checkmark & Mixed \\
\texttt{ext1}
  & $\times$ & \checkmark & \checkmark & Mixed \\
\texttt{lift}
  & $\times$ & $\times$ & $\times$ & Always fails \\
\texttt{measurability}
  & $\times$ & $\times$ & $\times$ & Always fails \\
\texttt{any\_goals}
  & \checkmark & \checkmark & $\times$ & Mixed \\
\texttt{by\_cases}
  & \checkmark & \checkmark & $\times$ & Mixed \\
\texttt{congrm}
  & $\times$ & $\times$ & $\times$ & Always fails \\
\bottomrule
\end{tabular}
\end{table}

\begin{finding}[The type-gated / architecture dichotomy]%
\label{find:dichotomy}
Schemas partition into two groups:
\begin{itemize}
\item \textbf{Type-gated} (\texttt{lift}, \texttt{measurability},
  \texttt{congrm}): require specific type-theoretic structures
  (CanLift instances, $\sigma$-algebra predicates, composite
  operator skeletons) that exist only in measure-theoretic areas.
  These fail in \emph{every} non-measure-theoretic target---0/9
  transfers across three areas.
\item \textbf{Architecture-sensitive}
  (\texttt{filter\_upwards}, \texttt{congr\,with}, \texttt{ext1},
  \texttt{any\_goals}, \texttt{by\_cases}): operate on proof
  structure (decomposition, case-split, extensionality) and
  transfer when the target has compatible goal shapes.  These
  succeed 8/15 times, but each schema's success depends on
  target-specific compatibility.
\end{itemize}
\end{finding}

The type-gated group is now confirmed with three independent
negative results per schema.  Future Yanasse runs should skip
\texttt{lift}, \texttt{measurability}, and \texttt{congrm}
for any non-measure-theoretic target, saving ${\sim}30$\% of
compute.

\paragraph{Why condensed is the hardest target.}
Condensed mathematics achieves a 20\% success rate versus 40\%
(RepTheory and SetTheory).  The difference is not
accidental: condensed proofs are \emph{categorically sequential}
(one goal at a time, handled by categorical lemma application),
while both other targets have goals with algebraic or combinatorial
structure amenable to case-splitting (\texttt{by\_cases}),
broadcasting (\texttt{any\_goals}), or conditional unfolding.

\begin{itemize}
\item \texttt{by\_cases} succeeds in RepTheory (morphism triviality
  split) and SetTheory (infinite/finite dichotomy) but fails in
  Condensed: no conditional definitions (\texttt{ite}/\texttt{dite})
  exist to unfold, and no structural dichotomies create two-branch
  proof architectures.
\item \texttt{any\_goals} succeeds in RepTheory (heterogeneous
  sub-goals after decomposition) and SetTheory (arithmetic
  side-conditions) but fails in Condensed: proofs never have
  more than one open goal simultaneously.
\item \texttt{filter\_upwards} fails in both SetTheory and
  Condensed, but for different reasons: SetTheory has no
  \texttt{Filter} API (ordinal arithmetic is algebraic);
  Condensed was \emph{designed} to replace filters with
  sheaf conditions.
\end{itemize}

%======================================================================
\section{The (Head, Modifier) Decomposition: Three-Area Evidence}%
\label{sec:finding}
%======================================================================

The (head, modifier) decomposition introduced in Part~1 receives
its strongest evidence from the three-area comparison of
\texttt{congr with}:

\begin{finding}[Full schema transfer]\label{find:full_transfer}
\texttt{congr with x} is the only schema to achieve full head +
modifier transfer in any target.  It does so only in
Condensed---the area where the target goal happens to be an
\emph{equality of constructor-shaped terms in a concrete category},
matching exactly the goal structure that \texttt{congr} expects.
In Representation Theory, only the modifier (\texttt{with})
transferred (the head required adaptation because the goal was
$\leq$, not $=$).  In Set Theory, both failed at rank-1
(Cardinal has no extensionality theorem), though a genuine transfer
existed at rank-3.
\end{finding}

This finding refines the (head, modifier) decomposition: whether
the head transfers depends on \emph{goal shape} (equality of
structured terms), while whether the modifier transfers depends on
\emph{proof structure} (whether per-component reduction is useful).
In the most favorable case---goal shape matches \emph{and}
per-component reduction is useful---the entire schema transfers
and may even \emph{improve} the existing proof.

%======================================================================
\section{Failures and Their Diagnoses}\label{sec:failures}
%======================================================================

The 8 failures share a common meta-diagnosis: condensed
mathematics was designed to \emph{eliminate} the structures that
probability-theory tactics operate on.

\begin{description}
\item[Item 20: \texttt{filter\_upwards}.]
  Requires \texttt{Filter.Eventually} goals.  Condensed mathematics
  has zero uses of the \texttt{Filter} API across all 33 Mathlib
  files.  The top-1 target (\texttt{isDiscrete\_tfae}) proves
  equivalence of 7 characterizations of discreteness---all
  categorical (adjunction counit isomorphisms, essential image
  membership, colimit characterizations).  This is a paradigm
  incompatibility: condensed math was designed to replace
  filter-based topology with sheaf conditions.

\item[Item 23: \texttt{all\_goals fun\_prop}.]
  Requires multiple simultaneous \texttt{Measurable}/\texttt{Continuous}
  goals.  The target (\texttt{isDiscrete\_tfae}) is a serial
  \texttt{tfae\_have} chain: never more than one goal open.
  Only 1~use of \texttt{fun\_prop} exists in all of
  \texttt{Mathlib/Condensed/}.

\item[Item 24: \texttt{lift}.]
  Requires \texttt{CanLift} instances (e.g.,
  $\mathbb{R}_{\geq 0}^\infty \to \mathbb{R}_{\geq 0}$).
  Condensed has zero \texttt{CanLift} instances, zero
  \texttt{WithTop}/\texttt{ENat}/\texttt{ENNReal} types.
  Replicates items~5 (RepTheory) and~14 (SetTheory).

\item[Item 25: \texttt{any\_goals}.]
  Requires multi-goal batching structure.  The target
  (\texttt{left\_triangle\_components}) has a single sequential
  goal at every proof step.  Source dispatches 4 symmetric
  bilinear-form sub-goals; target has no algebraic branching.
  Contrast with RepTheory (item~6, success: genuine multi-goal)
  and SetTheory (item~15, success: arithmetic side-conditions).

\item[Item 26: \texttt{measurability}.]
  A $\sigma$-algebra closure tactic (\texttt{aesop} with
  \texttt{@[measurability]} rule set).  Zero measurability goals
  exist in condensed proofs.  Replicates items~7 (RepTheory)
  and~16 (SetTheory).

\item[Item 27: \texttt{by\_cases}.]
  Requires conditional definitions (\texttt{ite}/\texttt{dite})
  to unfold.  Condensed has zero conditional definitions across
  all 33 files.  The target (\texttt{counit} naturality) is
  definitional: \texttt{ext; rfl}.  Contrast with SetTheory
  (item~17, success: infinite/finite dichotomy) and RepTheory
  (item~8, limited success: degenerate case split).

\item[Item 28: \texttt{congrm}.]
  Requires both sides of an equality to share a composite operator
  skeleton ($\int$, $\prod$, $\in$).  Target goal is an adjunction
  triangle identity ($f \gg g = \mathbb{1}$)---fundamentally
  asymmetric (composition vs.\ identity).  Replicates items~9
  (RepTheory) and~18 (SetTheory).

\item[Item 29: \texttt{split\_ifs}.]
  Requires \texttt{ite}/\texttt{dite} subexpressions in the goal.
  Zero \texttt{if-then-else} constructs exist in any condensed
  definition or proof across all 33 files.  Probability uses
  \texttt{split\_ifs} for indicator functions, density processes,
  and truncations---all domain-locked to areas with conditional
  definitions.
\end{description}

\paragraph{The common root cause.}
Six of the eight failures
(\texttt{filter\_upwards}, \texttt{all\_goals fun\_prop},
\texttt{lift}, \texttt{measurability}, \texttt{by\_cases},
\texttt{split\_ifs}) fail because the tactic requires a
\emph{type-theoretic precondition} that condensed mathematics
\emph{by design} never satisfies:

\medskip
\begin{center}\small
\begin{tabular}{@{}ll@{}}
\toprule
Schema & Required structure absent in Condensed \\
\midrule
\texttt{filter\_upwards} & \texttt{Filter}/\texttt{Filter.Eventually} \\
\texttt{all\_goals fun\_prop} & \texttt{Measurable}/\texttt{Continuous} goals \\
\texttt{lift} & \texttt{CanLift} instances \\
\texttt{measurability} & \texttt{@[measurability]}-tagged lemmas \\
\texttt{by\_cases} & Conditional definitions (\texttt{dite}) \\
\texttt{split\_ifs} & \texttt{ite}/\texttt{dite} subexpressions \\
\bottomrule
\end{tabular}
\end{center}

\medskip
\noindent The remaining two (\texttt{any\_goals}, \texttt{congrm})
fail for structural reasons: \texttt{any\_goals} requires
multi-goal proof states that condensed proofs' serial structure
never produces, and \texttt{congrm} requires symmetric composite
expressions that categorical morphism equalities lack.

%======================================================================
\section{Discussion}\label{sec:discussion}
%======================================================================

\subsection{The 20\% success rate}

The low success rate is informative, not disappointing.  It
confirms a prediction: when the target area was \emph{designed}
to eliminate the source area's mathematical vocabulary, only
schemas that operate at the level of proof \emph{architecture}
(decomposition into components) rather than proof
\emph{vocabulary} (filters, $\sigma$-algebras, conditional
definitions) can transfer.

The 20\% rate is consistent with the taxonomy from Part~1:
of the 10 schemas, 6 are type-gated (expected to fail), 2
are structurally sensitive (\texttt{any\_goals},
\texttt{congrm}---also expected to fail in this target), and
only 2 (\texttt{congr with}, \texttt{ext1}) are pure
decomposition patterns compatible with categorical
equality goals.

\subsection{Convergence to a single target theorem}

Both successful transfers and 6 of the 8 failures target
one of just two theorems:
\texttt{topCatAdjunction.left\_triangle\_components} and
\texttt{isDiscrete\_tfae}.  This convergence reflects the
small target pool (189 entries): the QAP matcher has limited
diversity to choose from, and a few high-scoring ``attractor''
entries dominate.  A larger condensed-mathematics formalization
would likely yield more diverse targets.

\subsection{The \texttt{set\_option} finding}

The most practically valuable result is P1's discovery that
\texttt{congr with x} can replace \texttt{ext} without the
\texttt{set\_option backward.isDefEq.respectTransparency false}
workaround.  This is not a theoretical observation: it produces
a compilable, \texttt{sorry}-free proof that is strictly simpler
than the Mathlib original.

The finding generalizes: any morphism equality in a concrete
category (TopCat, CompHaus, Profinite) where \texttt{ext} requires
the transparency override is a candidate for \texttt{congr with}
replacement.  The QAP-matched rank-2 analogue
(\texttt{Light/TopCatAdjunction.lean}) has the same proof
structure and likely benefits from the same replacement---a
testable prediction.

\subsection{Cost}

The full pipeline for this pair consumed:
\begin{itemize}
\item ${\sim}$5 minutes of GPU time per schema (MPS) for analogy
  matching ($8 \times 189$ comparisons)---the fastest pair in
  the series.
\item ${\sim}$6 hours of AI reasoning across 10 semantic adaptation
  sessions.
\item ${\sim}$\$80--120 in API cost.
\end{itemize}

%======================================================================
\section{Cumulative Results After Three Pairs}\label{sec:cumulative}
%======================================================================

\begin{table}[h]
\centering\small
\caption{Cumulative results across three Yanasse pairs.}
\label{tab:cumulative}
\begin{tabular}{@{}lccccl@{}}
\toprule
Pair & Schemas & Successes & Rate & Proofs & Pool \\
\midrule
Prob $\to$ RepTheory (Part~1)
  & 10 & 4 & 40\% & 4 & 672 \\
Prob $\to$ SetTheory (Part~2)
  & 10 & 4 & 40\% & 6 & 2{,}076 \\
Prob $\to$ Condensed (Part~3)
  & 10 & 2 & 20\% & 2 & 189 \\
\midrule
\textbf{Total} & \textbf{30} & \textbf{10} & \textbf{33\%}
  & \textbf{12} & --- \\
\bottomrule
\end{tabular}
\end{table}

Across 30 schema transfer attempts on 3 target areas, 10 schemas
succeeded (33\%), producing 12 Lean-verified new proofs.  The
``Proofs'' column exceeds ``Successes'' for Set Theory because
item~19 (\texttt{push}) yielded 3 independent proof variants.
The number of distinct target theorems receiving new proofs is~5:

\begin{enumerate}
\item \texttt{Coinvariants.le\_comap\_ker} (RepTheory, 2 proofs)
\item \texttt{$\varepsilon$ToSingle$_0$\_comp\_eq} (RepTheory, 2 proofs)
\item \texttt{mul\_eq\_left\_iff} (SetTheory, 3 schemas $\to$ 3 proofs)
\item \texttt{omega0\_lt\_preOmega\_iff} (SetTheory, 1 schema $\to$ 3 variants)
\item \texttt{left\_triangle\_components} (Condensed, 2 proofs)
\end{enumerate}

The most valuable individual result across all three papers is
the \texttt{congr with x} proof of
\texttt{left\_triangle\_components}, which is the only transfer
that produces a \emph{strictly better} proof than the Mathlib
original.

%======================================================================
\section{Conclusion}\label{sec:conclusion}
%======================================================================

The Probability $\to$ Condensed pair achieves a 20\% transfer
success rate, the lowest in the series, confirming that condensed
mathematics' deliberate elimination of measure-theoretic vocabulary
creates a hard barrier for most probability-theory schemas.

The positive results are correspondingly more valuable.  The
transferred \texttt{congr with x} tactic produces a proof that is
\emph{strictly better} than the Mathlib original---the only such
result in the series.  The finding that congruence-based
decomposition can replace extensionality-based decomposition
while avoiding transparency issues is a concrete, testable
insight that probability theory's ``go pointwise'' convention
contributes to condensed mathematics.

Cross-pair comparison across three target areas now confirms
the type-gated / architecture dichotomy with statistical
robustness: 0/9 transfers for type-gated schemas, 8/15 for
architecture-sensitive schemas, with target-specific compatibility
determining the outcome.  The permanently failing schemas
(\texttt{lift}, \texttt{measurability}, \texttt{congrm}) can be
pruned from future non-measure-theoretic runs.

Part~4 will apply the method to \texttt{Probability} $\to$
\texttt{FieldTheory}, the current target of the running Ralph Loop.

\section*{Acknowledgments}

Named after Horacio Hideki Yanasse.  The Deep Vision matching engine
was originally developed for chess cognition research; its application
to formal mathematics was enabled by the domain-independent design of
\texttt{deep\_vision\_lib.py}.  All computations were performed on a
MacBook Air (Apple M-series, 32\,GB unified memory) using PyTorch~2.8
on the MPS backend.

\begin{thebibliography}{99}

\bibitem{yanasse_part1}
A.~Linhares,
``Yanasse: Finding new proofs from Deep Vision's analogies --- Part~1:
Probability Theory $\to$ Representation Theory,''
ARGO LABS Report for Distribution, April 2026.

\bibitem{yanasse_part2}
A.~Linhares,
``Yanasse: Finding new proofs from Deep Vision's analogies --- Part~2:
Probability Theory $\to$ Set Theory,''
ARGO LABS Report for Distribution, April 2026.

\bibitem{linhares2024deepvision}
A.~Linhares,
``Deep Vision,''
ARGO LABS Internal Report, 2026.

\bibitem{hofstadter_fluid_1995}
D.~R.~Hofstadter and the Fluid Analogies Research Group, \emph{Fluid Concepts \& Creative Analogies: Computer Models of the Fundamental Mechanisms of Thought}. New York: Basic Books, 1995.

\bibitem{hofstadter_go_1999}
D.~R.~Hofstadter, \emph{G\"odel, Escher, Bach: An Eternal Golden Braid}, 20th anniversary ed. New York: Basic Books, 1999.

\bibitem{hofstadter_surfaces_2013}
D.~R.~Hofstadter and E.~Sander, \emph{Surfaces and Essences: Analogy as the Fuel and Fire of Thinking}. New York: Basic Books, 2013.

\bibitem{linhares_active_2005}
A.~Linhares, ``An active symbols theory of chess intuition,'' \emph{Minds and Machines}, vol.~15, pp.~131--151, 2005.

\bibitem{linhares_understanding_2007}
A.~Linhares and P.~Brum, ``Understanding our understanding of strategic scenarios: What role do chunks play?'' \emph{Cognitive Science}, vol.~31, no.~6, pp.~989--1007, 2007.

\bibitem{linhares_questioning_2010}
A.~Linhares and A.~E.~T.~A.~Freitas, ``Questioning Chase and Simon's (1973) `Perception in Chess','' \emph{New Ideas in Psychology}, vol.~28, no.~1, pp.~64--78, 2010.

\bibitem{linhares_entanglement_2012}
A.~Linhares, A.~E.~T.~A.~Freitas, A.~Mendes, and J.~S.~Silva, ``Entanglement of perception and reasoning in the combinatorial game of chess,'' \emph{Cognitive Systems Research}, vol.~13, no.~1, pp.~72--86, 2012.

\bibitem{linhares_what_2013}
A.~Linhares and D.~M.~Chada, ``What is the nature of the mind's pattern-recognition process?'' \emph{New Ideas in Psychology}, vol.~31, no.~2, pp.~108--121, 2013.

\bibitem{linhares_emergence_2014}
A.~Linhares, ``The emergence of choice: Decision-making and strategic thinking through analogies,'' \emph{Information Sciences}, vol.~259, pp.~36--56, 2014.

\bibitem{linhares_deep_vision}
A.~Linhares, ``Deep Vision: Seeing the essence of proofs through analogy,'' ARGO LABS Internal Report, 2026.

\bibitem{linhares_wolstenholme_2026}
A.~Linhares,
``Deep Vision: A formal proof of Wolstenholme's theorem in Lean~4,''
ARGO LABS Report for Distribution, 2026.

\bibitem{linhares_sketch_2026}
A.~Linhares,
``A sketch of Deep Vision's study of mathematics,''
ARGO LABS Report for Distribution, 2026.

\bibitem{scholze_condensed_2019}
P.~Scholze,
``Lectures on condensed mathematics,''
Lecture notes, University of Bonn, 2019.
Available at \url{https://www.math.uni-bonn.de/people/scholze/Condensed.pdf}.

\bibitem{clausen_scholze_complex_2022}
D.~Clausen and P.~Scholze,
``Condensed mathematics and complex geometry,''
Lecture notes, 2022.
Available at \url{https://people.mpim-bonn.mpg.de/scholze/Complex.pdf}.

\bibitem{mac_lane_categories_1998}
S.~Mac~Lane,
\emph{Categories for the Working Mathematician},
2nd ed., Graduate Texts in Mathematics, vol.~5.
New York: Springer-Verlag, 1998.

\bibitem{asgeirsson_condensed_2024}
D.~\'Asgeirsson, R.~Brasca, N.~Castellana, F.~Filipazzi,
E.~Kudryashov, P.~Lezama, and A.~Topaz,
``Categorical foundations of formalized condensed mathematics,''
\emph{J.~Symbolic Logic}, 2024.
\href{https://arxiv.org/abs/2407.12840}{arXiv:2407.12840}.

\bibitem{asgeirsson_discrete_2024}
D.~\'Asgeirsson,
``A formal characterization of discrete condensed objects,''
2024.
\href{https://arxiv.org/abs/2410.17847}{arXiv:2410.17847}.

\bibitem{scholze_liquid_2020}
P.~Scholze,
``Liquid tensor experiment,''
Blog post on Xena Project, December 2020.

\bibitem{commelin_lte_2022}
J.~Commelin, A.~Topaz, et al.,
``Completion of the Liquid Tensor Experiment,''
Lean Community Blog, 2022.

\end{thebibliography}

%======================================================================
% APPENDIX: Full Lean proofs
%======================================================================
\newpage
\appendix

\section{Lean Code: New Proofs}\label{app:lean}

\subsection{Item 21: \texttt{left\_triangle\_components} via
  \texttt{congr with x} (no \texttt{set\_option})}

\begin{lstlisting}
import Mathlib.Condensed.TopCatAdjunction
import Mathlib.Topology.Category.CompactlyGenerated

open Condensed CondensedSet CategoryTheory CompHaus

universe u

-- Transferred proof: congr with x (from Probability)
-- NOTE: No set_option needed!
example (Y : CondensedSet.{u}) :
    condensedSetToTopCat.map
      (topCatAdjunctionUnit Y) (*$\gg$*)
      topCatAdjunctionCounit
        (condensedSetToTopCat.obj Y) =
          (*$\mathbb{1}$*) _ := by
  congr with x
  change Y.obj.map ((*$\mathbb{1}$*) _) x = x
  simp
\end{lstlisting}

\newpage
\subsection{Item 22: \texttt{left\_triangle\_components} via
  \texttt{ext1 x} (three variants)}

\begin{lstlisting}
import Mathlib

open CategoryTheory Condensed CondensedSet CompHaus

universe u

-- Variant A: ext1 x + change + simp
set_option
    backward.isDefEq.respectTransparency false in
theorem left_triangle_A (Y : CondensedSet.{u}) :
    condensedSetToTopCat.map
      (topCatAdjunctionUnit Y) (*$\gg$*)
      topCatAdjunctionCounit Y.toTopCat =
        (*$\mathbb{1}$*) _ := by
  ext1 x
  change Y.obj.map ((*$\mathbb{1}$*) _) x = _
  simp

-- Variant B: ext1 x + change + simp only
set_option
    backward.isDefEq.respectTransparency false in
theorem left_triangle_B (Y : CondensedSet.{u}) :
    condensedSetToTopCat.map
      (topCatAdjunctionUnit Y) (*$\gg$*)
      topCatAdjunctionCounit Y.toTopCat =
        (*$\mathbb{1}$*) _ := by
  ext1 x
  change Y.obj.map ((*$\mathbb{1}$*) _) x = _
  simp only [FunctorToTypes.map_id_apply,
    condensedSetToTopCat_obj_carrier,
    TopCat.hom_id, ContinuousMap.id_apply]

-- Variant C: ext1 x + change + rw + rfl
set_option
    backward.isDefEq.respectTransparency false in
theorem left_triangle_C (Y : CondensedSet.{u}) :
    condensedSetToTopCat.map
      (topCatAdjunctionUnit Y) (*$\gg$*)
      topCatAdjunctionCounit Y.toTopCat =
        (*$\mathbb{1}$*) _ := by
  ext1 x
  change Y.obj.map ((*$\mathbb{1}$*) _) x = _
  rw [FunctorToTypes.map_id_apply]
  rfl
\end{lstlisting}

\end{document}
